\documentclass[12pt,a4paper,twoside]{article}
\usepackage[T1]{fontenc}
\usepackage[utf8]{inputenc}
\begin{document}

{\raggedleft Voici un paragraphe au fer à
droite où l’alignement ne se fait plus
qu’au niveau du bord droit. La syntaxe
anglaise peut être traduite par
\emph{déchiqueté~à~gauche}.\par}
{\raggedright Un paragraphe au fer à gauche.
La syntaxe anglaise est évidemment
cohérente en remplaçant gauche~par~droite.\par}
{\centering Paragraphe centré. Seules les
lignes relativement courtes ont un intérêt
sinon, on ne voit pas très bien le
caractère centré de la chose !
\par}
\end{document}